% GENERATED FILE. Edit canonical lessons, then run npm run generate:books.
% canonical-source-hash: fnv1a64:c7e4e5081e3ed196
% canonical-lessons: SA-C62-fifth, SA-C62-seventh, SA-C62-eighth, SA-C62-ninth, SA-C62-tenth, SA-R62-the-ending

\chapter{The Ending That Counts}
\label{ch:sa-ordinals}
\hlchaptermodality{62}

\begin{chapteropening}
\textbf{By the end of this chapter:} \emph{I can make a Sanskrit ordinal from a number by putting -ma on it, and I can name the one number between five and ten where that does not work.}

Complete SA-R62-the-ending and run pañcamaḥ to daśamaḥ, then name the number the chapter walked past and why.
\end{chapteropening}

\section[pañcamaḥ]{\sk{पञ्चमः} (pañcamaḥ) — fifth}

\label{lesson:SA-C62-fifth}



\begin{quote}
A second kind of number word begins here, so the recalls come first, at the four growing distances this book measures.

\begin{itemize}
\raggedright
  \item \emph{Recall:} read \textbf{\sk{ऋ}} — \textbf{R1}, one lesson back
  \item \emph{Recall:} say \emph{bhāratīyaḥ} — \textbf{R2}, five lessons back
  \item \emph{Recall:} close a quoted thought with \emph{iti} — \textbf{R3}, twenty lessons back
  \item \emph{Recall:} say \emph{sahasā} — \textbf{R4}, eighty lessons back
\end{itemize}
\end{quote}

\subsection*{You'll want to know: \sk{पञ्चमः}}
\textbf{\sk{पञ्चमः}} (\emph{pañcamaḥ}) — \textquotedblleft{}fifth\textquotedblright{}.

You have counted \textbf{\sk{एक} \sk{द्व} \sk{त्रि} \sk{चतुर्} \sk{पञ्च}} since chapter six. Those words answer \textbf{how many}. This one answers \textbf{which one in the line}, and Sanskrit keeps a second set of words for that job.

An ordinal is an ordinary adjective, so it wears the endings you already know: \textbf{\sk{पञ्चमः}} for a masculine thing, \textbf{\sk{पञ्चमम्}} for a neuter one.

\begin{cousinweb}
\textbf{\sk{पञ्च}} + \textbf{\sk{म}}.

The number is not worn down and not replaced. It stands there whole, and a single syllable goes on the end.

Macdonell's dictionary prints the word with the join marked — \emph{pañca-má} — and it will print four more the same way before this chapter is out.

\textbf{This lesson does not begin at \emph{first}.} It begins at \emph{fifth} because \textbf{-\sk{म}} is the ending that WORKS, and a rule has to be visible before anything can be seen departing from it. The words for \emph{first}, \emph{second} and \emph{third} all depart from it, and they wait.
\end{cousinweb}

\begin{grammarlens}[title={What you've built}]
One ordinal, and the ending that will make the next four for free.
\end{grammarlens}

\subsection*{Your turn}
Say these aloud:

\begin{itemize}
\raggedright
  \item \emph{pañca}, then \emph{pañcamaḥ}
  \item which of the two answers \emph{how many} — \emph{\textbf{pañca}}
  \item the ending that turned one into the other — \emph{\textbf{-ma}}
\end{itemize}

\begin{tcolorbox}[breakable,colback=teal!4,colframe=teal!35!black,title={Before you move on}]
What is \textbf{\sk{पञ्च}} plus \textbf{-\sk{म}}? (\textbf{\sk{पञ्चमः}}.) Why did this chapter not begin at \emph{first}? (\textbf{Because the rule has to be seen before a departure from it can be read as one.})

Source: Macdonell, \emph{A Practical Sanskrit Dictionary}, p. 149 (DSAL).
\end{tcolorbox}
\section[saptamaḥ]{\sk{सप्तमः} (saptamaḥ) — seventh}

\label{lesson:SA-C62-seventh}



\begin{quote}
Four recalls, at four distances.

\begin{itemize}
\raggedright
  \item \emph{Recall:} say the ending that makes an ordinal — \textbf{R1}, one lesson back
  \item \emph{Recall:} say \emph{Bhāratāt} and name the ending that means \textquotedblleft{}from\textquotedblright{} — \textbf{R2}, five lessons back
  \item \emph{Recall:} say \emph{iti} — \textbf{R3}, twenty-two lessons back
  \item \emph{Recall:} say \emph{prāyaḥ} — \textbf{R4}, eighty lessons back
\end{itemize}
\end{quote}

\subsection*{You'll want to know: \sk{सप्तमः}}
\textbf{\sk{सप्तमः}} (\emph{saptamaḥ}) — \textquotedblleft{}seventh\textquotedblright{}.

Build it before you read it: you have \textbf{\sk{सप्त}}, and you have \textbf{-\sk{म}}.

\textbf{\sk{सप्त}} + \textbf{\sk{म}}. Nothing else changed.

\begin{cousinweb}
Chapter nineteen showed you that \textbf{\sk{सप्त}}, Latin \emph{septem} and English \emph{seven} are the same word travelling three roads. The ordinal carries the same news one step further.

Latin's seventh is \textbf{septimus}. Look at the two endings side by side: Sanskrit \textbf{-\sk{म}}, Latin \textbf{-mus}. That is not a coincidence and it is not a borrowing — it is one inherited ordinal ending, worn slightly differently in two mouths.

English went the other way and used a different suffix for all of them, which is why English says \emph{seven-th} where Latin says \emph{septi-mus}. Both endings are old. Sanskrit kept both, and the rest of this chapter and the next is the story of which numbers got which.
\end{cousinweb}

\begin{grammarlens}[title={What you've built}]
A second ordinal, made rather than learned, and the first sign that the ending is older than Sanskrit.
\end{grammarlens}

\subsection*{Your turn}
Say these aloud:

\begin{itemize}
\raggedright
  \item \emph{sapta}, then \emph{saptamaḥ}
  \item \emph{pañcamaḥ}, then \emph{saptamaḥ}
  \item the Latin word with the same ending — \emph{septimus}
\end{itemize}

\begin{tcolorbox}[breakable,colback=teal!4,colframe=teal!35!black,title={Before you move on}]
What is \textbf{\sk{सप्त}} plus \textbf{-\sk{म}}? (\textbf{\sk{सप्तमः}}.) Which Latin word shows the same ending? (\emph{Septimus}.)

Source: Macdonell, \emph{A Practical Sanskrit Dictionary}, p. 335 (DSAL).
\end{tcolorbox}
\section[aṣṭamaḥ]{\sk{अष्टमः} (aṣṭamaḥ) — eighth}

\label{lesson:SA-C62-eighth}



\begin{quote}
Four recalls, at four distances.

\begin{itemize}
\raggedright
  \item \emph{Recall:} say \emph{saptamaḥ} — \textbf{R1}, one lesson back
  \item \emph{Recall:} draw the \textbf{\sk{।}} and say what it marks — \textbf{R2}, five lessons back
  \item \emph{Recall:} say \emph{gatvā} and name what the \textbf{-\sk{त्वा}} does — \textbf{R3}, twenty lessons back
  \item \emph{Recall:} say \emph{kiñcit} — \textbf{R4}, eighty lessons back
\end{itemize}
\end{quote}

\subsection*{You'll want to know: \sk{अष्टमः}}
\textbf{\sk{अष्टमः}} (\emph{aṣṭamaḥ}) — \textquotedblleft{}eighth\textquotedblright{}.

\textbf{\sk{अष्ट}} + \textbf{\sk{म}}. Third time, same move.

\begin{cousinweb}
Eight is where the cousins stop agreeing. Latin's eighth is \textbf{octāvus} — not \emph{octimus}, not \emph{octus}, but a shape that matches nothing else in its own row. English says \emph{eighth}, keeping its own \textbf{-th}.

Sanskrit does nothing unusual at all here. \textbf{\sk{अष्ट}} takes \textbf{-\sk{म}} exactly as \textbf{\sk{सप्त}} did.

That is worth noticing for a reason beyond the word: it tells you that the \textbf{-\sk{म}} ending is not merely common in Sanskrit but PRODUCTIVE — the ending a speaker reaches for by default, rather than one of several stored shapes. You will meet the one number in this range where that is not true, and it will stand out.
\end{cousinweb}

\begin{grammarlens}[title={What you've built}]
A third ordinal for free, and a reason to expect the fourth.
\end{grammarlens}

\subsection*{Your turn}
Say these aloud:

\begin{itemize}
\raggedright
  \item \emph{aṣṭa}, then \emph{aṣṭamaḥ}
  \item \emph{pañcamaḥ saptamaḥ aṣṭamaḥ}
  \item which language broke step at eight — \textbf{Latin}
\end{itemize}

\begin{tcolorbox}[breakable,colback=teal!4,colframe=teal!35!black,title={Before you move on}]
What is \textbf{\sk{अष्ट}} plus \textbf{-\sk{म}}? (\textbf{\sk{अष्टमः}}.) Did Sanskrit do anything special at eight? (\textbf{No.})

Source: Macdonell, \emph{A Practical Sanskrit Dictionary}, p. 33 (DSAL).
\end{tcolorbox}
\section[navamaḥ]{\sk{नवमः} (navamaḥ) — ninth}

\label{lesson:SA-C62-ninth}



\begin{quote}
Four recalls, at four distances.

\begin{itemize}
\raggedright
  \item \emph{Recall:} say \emph{aṣṭamaḥ} — \textbf{R1}, one lesson back
  \item \emph{Recall:} say \emph{ṛtuḥ} — \textbf{R2}, five lessons back
  \item \emph{Recall:} say \emph{gantum} and name what the \textbf{-\sk{तुम्}} does — \textbf{R3}, twenty lessons back
  \item \emph{Recall:} say \emph{bahu} — \textbf{R4}, eighty lessons back
\end{itemize}
\end{quote}

\subsection*{You'll want to know: \sk{नवमः}}
\textbf{\sk{नवमः}} (\emph{navamaḥ}) — \textquotedblleft{}ninth\textquotedblright{}.

\textbf{\sk{नव}} + \textbf{\sk{म}}. Chapter nineteen taught the number as \textbf{\sk{नवन्}}; the final \textbf{\sk{न्}} drops before the ending and \textbf{\sk{नव}} is what the ordinal is built on.

\begin{cousinweb}
Now a warning, because two Sanskrit words meet here and look identical.

\begin{quote}
\textbf{\sk{नव}} — \emph{nava} — \textquotedblleft{}\textbf{new}\textquotedblright{} (chapter fifteen)
\end{quote}

\begin{quote}
\textbf{\sk{नव}} — \emph{nava} — \textquotedblleft{}\textbf{nine}\textquotedblright{}, the stem under \textbf{\sk{नवन्}} (chapter nineteen)
\end{quote}

They are not one word with two senses. They are two different words that happened to wear down to the same shape, and Sanskrit lives with the collision.

\textbf{\sk{नवमः}} is built on the SECOND of them. It means \emph{ninth}, never \emph{newest}.

Latin has the same pair and the same trouble — \emph{novus} \textquotedblleft{}new\textquotedblright{} beside \emph{novem} \textquotedblleft{}nine\textquotedblright{} — which is a good reason to believe the collision is inherited rather than invented in India.
\end{cousinweb}

\begin{grammarlens}[title={What you've built}]
A fourth ordinal, and a pair of look-alikes told apart.
\end{grammarlens}

\subsection*{Your turn}
Say these aloud:

\begin{itemize}
\raggedright
  \item \emph{navan}, then \emph{navamaḥ}
  \item the other \textbf{\sk{नव}}, the one that means \emph{new}
  \item \emph{saptamaḥ aṣṭamaḥ navamaḥ}
\end{itemize}

\begin{tcolorbox}[breakable,colback=teal!4,colframe=teal!35!black,title={Before you move on}]
Which \textbf{\sk{नव}} is \textbf{\sk{नवमः}} built on? (\textbf{The one that means nine.}) What does the other one mean? (\textbf{New.})

Source: Macdonell, \emph{A Practical Sanskrit Dictionary}, p. 137 (DSAL).
\end{tcolorbox}
\section[daśamaḥ]{\sk{दशमः} (daśamaḥ) — tenth}

\label{lesson:SA-C62-tenth}



\begin{quote}
Four recalls, at four distances.

\begin{itemize}
\raggedright
  \item \emph{Recall:} say \emph{navamaḥ} — \textbf{R1}, one lesson back
  \item \emph{Recall:} read \textbf{\sk{ऋ}} — \textbf{R2}, five lessons back
  \item \emph{Recall:} say \emph{patiḥ} — \textbf{R3}, forty lessons back
  \item \emph{Recall:} say \emph{ātithyam} — \textbf{R4}, eighty lessons back
\end{itemize}
\end{quote}

\subsection*{You'll want to know: \sk{दशमः}}
\textbf{\sk{दशमः}} (\emph{daśamaḥ}) — \textquotedblleft{}tenth\textquotedblright{}.

\textbf{\sk{दश}} + \textbf{\sk{म}}, and Latin's tenth is \textbf{decimus} — the same ending again, in the same place.

\begin{cousinweb}
Count what this chapter has done:

\noindent
\begin{tabularx}{\linewidth}{@{}>{\raggedright\arraybackslash}X>{\raggedright\arraybackslash}X@{}}
\toprule
\textbf{number} & \textbf{ordinal} \\
\midrule
\textbf{\sk{पञ्च}} & \textbf{\sk{पञ्चमः}} \\
\textbf{\sk{सप्त}} & \textbf{\sk{सप्तमः}} \\
\textbf{\sk{अष्ट}} & \textbf{\sk{अष्टमः}} \\
\textbf{\sk{नवन्}} & \textbf{\sk{नवमः}} \\
\textbf{\sk{दश}} & \textbf{\sk{दशमः}} \\
\bottomrule
\end{tabularx}

Five ordinals, one ending, nothing memorised.

Now look at the left-hand column and find what is missing. \textbf{\sk{षट्}} is not there. You have counted with it since chapter nineteen and this chapter has walked straight past it.

That was on purpose, and it is not because \emph{sixth} is rare. \textbf{\sk{षट्}} is the one number between five and ten whose ordinal does NOT take \textbf{-\sk{म}}, and the next chapter opens on it — because a hole is only interesting once the thing around it is solid.
\end{cousinweb}

\begin{grammarlens}[title={What you've built}]
Five ordinals from one ending, and one number deliberately left out.
\end{grammarlens}

\subsection*{Your turn}
Say these aloud:

\begin{itemize}
\raggedright
  \item \emph{daśa}, then \emph{daśamaḥ}
  \item the run — \emph{pañcamaḥ saptamaḥ aṣṭamaḥ navamaḥ daśamaḥ}
  \item the number this chapter skipped — \emph{\textbf{ṣaṭ}}
\end{itemize}

\begin{tcolorbox}[breakable,colback=teal!4,colframe=teal!35!black,title={Before you move on}]
How many ordinals did this chapter build from one ending? (\textbf{Five}.) Which number did it leave out, and why? (\emph{\textbf{\sk{षट्}}} — \textbf{because its ordinal does not take -\sk{म}}.)

Source: Macdonell, \emph{A Practical Sanskrit Dictionary}, p. 117 (DSAL).
\end{tcolorbox}
\section[pañcamaḥ · saptamaḥ · aṣṭamaḥ …]{Five ordinals, one ending, and a hole at six}

\label{lesson:SA-R62-the-ending}



\begin{quote}
Four recalls, at four distances, before the run.

\begin{itemize}
\raggedright
  \item \emph{Recall:} say \emph{daśamaḥ} — \textbf{R1}, one lesson back
  \item \emph{Recall:} say \emph{pañcamaḥ} — \textbf{R2}, five lessons back
  \item \emph{Recall:} say \emph{agacchat} and name the \textbf{\sk{अ}-} at the front — \textbf{R3}, twenty lessons back
  \item \emph{Recall:} say \emph{vinayaḥ} — \textbf{R4}, eighty lessons back
\end{itemize}
\end{quote}

\subsection*{Your turn}
Say these aloud:

\begin{itemize}
\raggedright
  \item the five in order — \emph{pañcamaḥ saptamaḥ aṣṭamaḥ navamaḥ daśamaḥ}
  \item each number with its ordinal — \emph{pañca pañcamaḥ}, \emph{sapta saptamaḥ}, and on
  \item how many of the five you had to be told — \textbf{none of them}
  \item the number missing from the run, and what it is — \emph{\textbf{ṣaṭ}}, \textbf{six}
\end{itemize}

\begin{grammarlens}[title={What you've built}]
Five ordinals for the price of one syllable, and one question left open.

The question is worth stating exactly, because the next chapter answers it and then keeps going. It is not \emph{\textquotedblleft{}why is sixth irregular?\textquotedblright{}} — that phrasing assumes the ending you learned is the only one Sanskrit has. It is:

\begin{quote}
\textbf{If -\sk{म} is not what sixth uses, what does it use, and is it alone?}
\end{quote}

It is not alone. There are two more endings in this system, and between them they account for every ordinal below five as well as the one this chapter walked past.
\end{grammarlens}

\begin{tcolorbox}[breakable,colback=teal!4,colframe=teal!35!black,title={Before you move on}]
What ending makes an ordinal out of a number? (\textbf{-\sk{म}}.) Which number between five and ten does not use it? (\emph{\textbf{\sk{षट्}}}.)
\end{tcolorbox}
